\chapter{力位混合、笛卡尔阻抗与全身控制（WBC/TSID-HQP）}
\label{ch:force-wbc}

% ════════════════════════════════════════════════════════════════
%  卷五《机器人最优控制》· 控制部「浮基/接触/力控」子部 P 章（收口）。
%  主线：机械臂侧力控（Raibert-Craig → Salisbury → Chiaverini → Hogan 阻抗）
%        → 浮基升维（12 维顶点力 / 二端口四任务 / 阻抗↔WBC 统一 KKT）。
%  与 ch:quad_wbc（腿足侧 WBC）、ch:control（导论）、O 章（操作空间）强互引去重。
% ════════════════════════════════════════════════════════════════

\cref{sec:ctrl-robot-impedance} 在控制导论里给出了机器人与环境交互的一个核心洞见：接触任务不应单独控制位置或力，而应控制末端的\emph{动态关系}——机械阻抗。那里给出的是一个目标阻抗方程与免力传感特例，作为入门级的“一张图景”。本章把这条线索展开成机器人力控的\emph{完整谱系}：从 1980 年代奠基的力位混合、主动刚度、并行力位，到 Hogan 阻抗的完整稳定性证明，再到浮动基座下的全身控制（WBC）与任务空间逆动力学（TSID）。

贯穿全章的有一条主梁——\textbf{“位置控制 $\to$ 力位混合 $\to$ 阻抗（含无源性）$\to$ WBC（多任务约束优化）”}。我们将看到这条线并非五种互不相干的技巧，而是同一物理问题（如何在接触下同时调控力与运动）在“约束越来越复杂、自由度越来越高”的方向上逐级推广。本章定位为这条线的\emph{机械臂侧入口}：固定基座串联臂的五大经典算法 + 阻抗的稳定性深化 + 面接触 $12$ 维顶点力 + “阻抗是 WBC 单任务特例”的统一桥。它与 \cref{ch:quad_wbc}（从腿足/浮基进入 WBC、点接触 $3$D 力、MPC 接力）互补：两章在“WBC $=$ 多任务约束优化”处会合并互引。零空间投影、伪逆最小范数证明、$42$ 维浮基 QP、HQP 字典序最小化等机制已在 \cref{ch:quad_wbc} 讲透，本章一律引用复用，绝不重写——把笔墨集中在那里没有的内容上。

\begin{remark}[本章的符号约定，全章一把尺]
\label{rmk:force-wbc-notation}
为与 \cref{ch:control,ch:quad_wbc} 不打架，本章统一约定：
\begin{itemize}
  \item \textbf{位置误差} $\mathbf{e}_x=\mathbf{x}_d-\mathbf{x}$；恢复力写 $+\mathbf{K}_d^x\mathbf{e}_x$。引用源里 $\tilde{\mathbf{x}}=\mathbf{x}-\mathbf{x}_d=-\mathbf{e}_x$ 时显式桥接。
  \item \textbf{姿态误差}用 \cref{ch:lie} 的右扰动对数映射 $\mathbf{e}_R=\mathrm{Log}(\mathbf{R}^\top\mathbf{R}_d)\in\mathbb{R}^3$（机体系），恢复律 $\boldsymbol{\omega}^*=\mathbf{K}_R\mathbf{e}_R+\cdots$，与 \cref{sec:ctrl-robot-osc} 逐字一致。
  \item \textbf{两种“选择矩阵”严格区分}：力位混合的\emph{任务选择矩阵} $\mathbf{S}_\sigma=\mathrm{diag}(s_i),\,s_i\in\{0,1\}$（$6\times6$）；浮基\emph{驱动选择矩阵} $\mathbf{S}=[\mathbf{0}\ \ \mathbf{I}_{n_j}]$（$n_j\times(6+n_j)$）。
  \item \textbf{两种伪逆严格区分}：动力学一致逆 $\bar{\mathbf{J}}=\mathbf{M}^{-1}\mathbf{J}^\top\boldsymbol{\Lambda}$（操作空间正解），与 Moore--Penrose 逆 $\mathbf{J}^+$（纯运动学最小范数，做零空间会“泄漏”到末端）。
  \item 操作空间惯量 $\boldsymbol{\Lambda}=(\mathbf{J}\mathbf{M}^{-1}\mathbf{J}^\top)^{-1}$、任务科氏 $\boldsymbol{\mu}$、任务重力 $\mathbf{p}$ 沿用 \cref{eq:ctrl-osc}；广义速度 $\boldsymbol{\nu}$、广义加速度 $\dot{\boldsymbol{\nu}}$、接触力 $\boldsymbol{\lambda}$、接触雅可比 $\mathbf{J}_c$ 沿用 \cref{eq:ctrl-floatbase}。
\end{itemize}
\end{remark}

% ════════════════════════════════════════════════════════════════
\section{力控的问题与历史}
\label{sec:fw-history}

\begin{insight}[力控发展的历史脊柱：一条“调什么”的演进主线]
\label{ins:fw-spine}
进入具体推导前，先把全部力控思想串成一条\emph{为何这样演进}的脊柱。\cref{fig:fw-timeline} 与下文“三幕历史”按\emph{时间/人物}排（谁、何年），此处则按\emph{“到底在调控什么”}排——这是理解一切力控方案的真正主线。三幕回答的是同一个递进问题：接触时既不能只控位置、又难以只控力（\cref{tab:fw-stability} 印证“只控力须知 $K_e$”），那么\textbf{该把什么当作被设计的对象}？
\begin{itemize}
  \item \textbf{第一幕——调“柔顺/刚度”（把环境顺应度本身做成可设计量）。}最早的答案：不去争力或位，而是让末端\emph{软}下来，用顺应吸收接触误差。先有 Whitney 的\emph{被动柔顺}（RCC，纯机械弹性，零传感）；再有 Salisbury 的\emph{主动刚度}\cite{salisbury1980active}——用软件把末端整形成一个虚拟弹簧 $\tilde{\mathbf{x}}=\mathbf{K}_x^{-1}\mathbf{f}_{\rm ext}$（\cref{sec:fw-stiffness}）。同期 Mason 给出任务侧的语言：每个方向上环境要么限位、要么限力（\emph{自然约束}），余者才可由控制器指定（\emph{人工约束}）\cite{mason1981compliance}——这套\emph{约束分析}是后两幕的公共地基。其软肋是纯弹簧无阻尼会\emph{弹跳}（\cref{ins:fw-stiffness-bounce}），逼出第二幕。
  \item \textbf{第二幕——调“力--运动关系”（Hogan 阻抗的统一框架）。}Hogan（1985）把第一幕的“静态刚度”升格为整条\emph{二阶动态关系}\cite{hogan1985impedance}：不分别控力或控位，而是规定末端的虚拟\emph{质量--弹簧--阻尼} $\mathbf{M}_d\ddot{\tilde{\mathbf{x}}}+\mathbf{D}_d\dot{\tilde{\mathbf{x}}}+\mathbf{K}_d\tilde{\mathbf{x}}=\mathbf{f}_{\rm ext}$（\cref{sec:fw-impedance}）。一个旋钮（$\mathbf{K}_d$）就把“硬位置控制”连续旋到“软柔顺”，且自带\emph{无源性}保证——对从泡沫到钢铁的全谱环境都稳定，无需知道 $K_e$（\cref{sec:fw-passivity}）。这就是\emph{间接力控}的精髓：力不被直接闭环，而由“关系 $+$ 运动”间接定。其对偶实现（测力出位移）即\emph{导纳}控制，二者是同一关系的两种因果取向。
  \item \textbf{第三幕——调“力位的子空间分解”（Raibert--Craig 混合 / 受约束力位控制）。}当任务几何\emph{已知}（擦平面、插轴），可比阻抗更进一步：借第一幕的约束分析，用\emph{选择矩阵} $\mathbf{S}_\sigma$ 把任务帧一刀切成“控位子空间”与“控力子空间”，各自显式闭环（\cref{sec:fw-raibert}）——这是\emph{直接力控}，力被真正闭环到设定值。Chiaverini 的\emph{并行力位}\cite{chiaverini1993parallel} 则用增益大小做隐式分解、免去硬切换（\cref{sec:fw-parallel}）。代价是须先知约束帧、且稳定性\emph{依赖} $K_e$（\cref{ins:fw-rc-cond}）——这反衬出第二幕无源性的可贵。
\end{itemize}
拿到这条脊柱再读下文：第一幕\cref{sec:fw-stiffness}、第二幕\cref{sec:fw-impedance,sec:fw-passivity}、第三幕\cref{sec:fw-raibert,sec:fw-parallel} 只是把每一幕“调什么”的承诺兑现成可编程的控制律；而\cref{sec:fw-unification} 的 WBC 不过是把“调力--运动关系”推广到浮基多任务的终点站。
\end{insight}

\paragraph{为什么接触时纯位置控制会“打架”。}
设末端在接触方向上的等效闭环刚度为 $K$，与刚度 $K_e$ 的环境相碰。一个微小的位置跟踪误差 $\delta x$ 会被放大成接触力 $f\approx \dfrac{K K_e}{K+K_e}\,\delta x$。工业位置环为追求跟踪精度，$K$ 常达 $10^4$--$10^6$ 量级；碰到钢铁（$K_e>10^6$）时，亚毫米误差即产生数百牛顿的力——要么硬碰硬出现力尖峰与跳变，要么把增益全程调软牺牲自由空间精度。\textbf{根本原因是：在受约束方向上，位置与力不能同时独立指定}（环境约束把二者耦合，\cref{ch:control} 的“接触力大就把位置环调软”误区即源于此）。出路有两条：在\emph{已知}接触几何上分管力与位置（力位混合），或显式设定力--运动\emph{关系}（阻抗）。本章把这两条路及其后继全部铺开。

\paragraph{三幕历史。}
机器人力控的思想史可分三幕（\cref{fig:fw-timeline}）。\textbf{第一幕：被动柔顺。}Whitney 等在 Draper 实验室于 1970 年代中期研发了\emph{遥心柔顺}（Remote Center Compliance, RCC）——一个纯机械的被动顺应装置（Drake 1977 的 MIT 博论给出方法，Whitney--Nevins 1979 系统阐述），用机械弹性吸收装配误差，不需任何传感或计算。\textbf{第二幕：主动力位分管。}Mason（1981）先把接触任务的几何抽象为\emph{自然约束}与\emph{人工约束}\cite{mason1981compliance}（约束分析必须先于控制器设计）；Raibert \& Craig（1981）据此提出\emph{力位混合控制}\cite{raibert1981hybrid}——这是被引数千次的奠基性工作，用选择矩阵把任务空间分成“控位”与“控力”两组子空间分别施控。\textbf{第三幕：阻抗统一。}Hogan（1985）的三部曲\cite{hogan1985impedance} 指出：与其分别控制力或位置，不如把末端整形为一个虚拟的\emph{质量--弹簧--阻尼}（机械阻抗），用一个参数从“硬”连续过渡到“软”，并自带与被动环境耦合的稳定性。其间，Salisbury（1980）的主动刚度\cite{salisbury1980active} 是阻抗的简化前身，De Schutter \& Van Brussel（1988）的\emph{任务帧/柔顺帧}\cite{deschutter1988compliant} 给了混合控制以坐标系语言，Chiaverini \& Sciavicco（1993）的\emph{并行力位}\cite{chiaverini1993parallel} 则免去了选择矩阵。本章 \cref{sec:fw-raibert,sec:fw-direct,sec:fw-stiffness,sec:fw-parallel,sec:fw-impedance} 依次展开这五条主线。

\begin{figure}[H]
\centering
\begin{tikzpicture}[>=Stealth, font=\small,
  ev/.style={align=center, font=\footnotesize}]
  \draw[->, thick, gray!60!black] (0,0) -- (13,0);
  % 每条数据自带：x / 文字锚点纵坐标 y / 锚点 anc / 茎端 sy（避免运行时条件判断）
  \foreach \x/\y/\anc/\sy/\name/\who in {%
    0.6/1.0/south/0.18/{1975--79}/{被动柔顺 RCC\\ Whitney/Drake},
    3.4/-1.0/north/-0.18/{1980}/{主动刚度\\ Salisbury},
    5.6/1.0/south/0.18/{1981}/{力位混合\\ Raibert--Craig\\（Mason 约束分析）},
    8.0/-1.0/north/-0.18/{1985}/{阻抗控制\\ Hogan},
    10.2/1.0/south/0.18/{1988}/{任务帧 / 耦合稳定\\ De Schutter; Colgate},
    12.2/-1.0/north/-0.18/{1993}/{并行力位\\ Chiaverini}}{%
    \fill[main!70!black] (\x,0) circle (2.2pt);
    \node[ev, anchor=\anc] at (\x,\y) {\textbf{\name}\\ \who};
    \draw[gray!50] (\x,0) -- (\x,\sy);
  }
  \node[anchor=west, font=\footnotesize\itshape] at (0,-1.95)
    {被动 $\longrightarrow$ 主动力位分管 $\longrightarrow$ 阻抗统一};
\end{tikzpicture}
\caption{机器人力控思想史三幕。被动柔顺（机械顺应）$\to$ 主动力位分管（选择矩阵 / 主动刚度）$\to$ 阻抗统一（整形力--运动关系，自带无源稳定性）。}
\label{fig:fw-timeline}
\end{figure}

% ════════════════════════════════════════════════════════════════
\section{Raibert--Craig 力位混合控制}
\label{sec:fw-raibert}

\paragraph{动机：力与位置分管不同方向。}
擦桌子、磨平面、插轴这类任务的共性是：在\emph{某些}笛卡尔方向上要控\emph{位置}（沿桌面运动），在\emph{另一些}与之正交的方向上要控\emph{力}（压住桌面）。Raibert--Craig 的思想朴素而有力——用一个对角的\textbf{任务选择矩阵}把六维任务空间一刀切开。

\begin{definition}[任务选择矩阵与约束帧]
\label{def:fw-selection}
在\emph{约束帧}（亦称任务帧 / 柔顺帧，De Schutter 1988\cite{deschutter1988compliant}）下定义对角矩阵
\begin{equation}
  \mathbf{S}_\sigma=\mathrm{diag}(s_1,\dots,s_6),\qquad s_i\in\{0,1\},
  \label{eq:fw-selection}
\end{equation}
$s_i=1$ 表示第 $i$ 个任务方向\emph{控位置}，$s_i=0$ 表示\emph{控力}；其补 $\bar{\mathbf{S}}_\sigma=\mathbf{I}-\mathbf{S}_\sigma$。当约束帧相对世界系有旋转 $\mathbf{R}_{cf}$ 时，世界系下的选择算子为 $\mathbf{S}_{\rm world}=\mathbf{T}_{cf}\,\mathbf{S}_\sigma\,\mathbf{T}_{cf}^\top$，其中 $\mathbf{T}_{cf}=\mathrm{blkdiag}(\mathbf{R}_{cf},\mathbf{R}_{cf})$ 把力与力矩两块同时旋转（旋转矩阵 $\mathbf{R}_{cf}\in\SO(3)$ 见 \cref{ch:rigid_body}）。
\end{definition}

\paragraph{完整控制律。}
分三步把任务描述映射到关节力矩。\emph{Step 1}（任务空间控制力）：控位子空间用 PD，控力子空间用 PI，
\begin{equation}
  \mathbf{F}_{\rm pos}=\mathbf{K}_p\mathbf{e}_x+\mathbf{K}_d\dot{\mathbf{e}}_x,\qquad
  \mathbf{F}_{\rm force}=\mathbf{K}_{fp}\mathbf{e}_f+\mathbf{K}_{fi}\!\int_0^t\!\mathbf{e}_f\,\mathrm{d}\tau,
  \label{eq:fw-rc-forces}
\end{equation}
$\mathbf{e}_x=\mathbf{x}_d-\mathbf{x}$、$\mathbf{e}_f=\mathbf{f}_d-\mathbf{f}$。\emph{Step 2}（选择矩阵组合）：$\mathbf{F}_{\rm task}=\mathbf{S}_\sigma\mathbf{F}_{\rm pos}+\bar{\mathbf{S}}_\sigma\mathbf{F}_{\rm force}$。\emph{Step 3}（对偶映射到关节，含重力/科氏补偿）：
\begin{equation}
  \boxed{\;\boldsymbol{\tau}=\mathbf{J}^\top\!\Big[\mathbf{S}_\sigma\big(\mathbf{K}_p\mathbf{e}_x+\mathbf{K}_d\dot{\mathbf{e}}_x\big)+\bar{\mathbf{S}}_\sigma\big(\mathbf{K}_{fp}\mathbf{e}_f+\mathbf{K}_{fi}\!\textstyle\int\mathbf{e}_f\big)\Big]+\mathbf{C}\dot{\mathbf{q}}+\mathbf{g}.\;}
  \label{eq:fw-rc-law}
\end{equation}

\begin{remark}[力控接口的量纲二义性——一个反复被踩的坑]
\label{rmk:fw-interface}
\cref{eq:fw-rc-forces} 的力 PI 输出究竟是什么，决定了增益的量纲，二者不可混用。\textbf{力接口（wrench）}：PI 直接输出任务空间 wrench，则 $\mathbf{K}_{fp}$ \emph{无量纲}、$\mathbf{K}_{fi}$ 量纲 $1/\mathrm{s}$，输出经 $\mathbf{J}^\top$ 映射到力矩，如 \cref{eq:fw-rc-law}。\textbf{位置偏置接口（导纳式）}：PI 输出位置偏置 $\Delta\mathbf{x}$ 送入内层位置/导纳环，则 $\mathbf{K}_{fp}$ 量纲 $\mathrm{m/N}$、$\mathbf{K}_{fi}$ 量纲 $\mathrm{m/(N\!\cdot\! s)}$，且\emph{不能}再直接并入 $\mathbf{J}^\top$。推导与实现必须\emph{先选定接口、再写单位与闭环模型}——把 $\mathrm{m/N}$ 增益塞进 $\mathbf{J}^\top$ wrench 律是常见错误。
\end{remark}

\paragraph{稳定性：依赖环境刚度。}
控位子空间（$s_i=1$）闭环为标准二阶系统 $\Lambda_{ii}\ddot{\tilde{x}}_i+K_{d,i}\dot{\tilde{x}}_i+K_{p,i}\tilde{x}_i=0$，$K_{p,i},K_{d,i}>0$ 即渐近稳定（标准 PD 结论，见 \cref{ch:lyapunov-stability}）。控力子空间（$s_i=0$）则要把力 PI 与弹性环境 $f=K_e\tilde{x}$ 联立（此处取\emph{位置偏置接口}的外环增益），闭环特征方程为\emph{三阶}：
\begin{equation}
  M_d s^3+D_{\rm eff}s^2+K_e(1+K_{fp})s+K_e K_{fi}=0,
  \label{eq:fw-rc-char}
\end{equation}
由 Routh--Hurwitz 判据（三阶系统稳定的充要条件为各系数同号且 $a_2a_1>a_3a_0$，见 \cref{ch:lyapunov-stability}），稳定要求
\begin{equation}
  D_{\rm eff}(1+K_{fp})>M_d K_{fi}.
  \label{eq:fw-rc-rh}
\end{equation}
注意 $K_e$ 同时出现在 \cref{eq:fw-rc-char} 的 $s^1$ 与 $s^0$ 系数中，作 $a_2a_1>a_3a_0$ 时\emph{恰好相消}——故此集总三阶模型给出的必要条件\emph{不显含} $K_e$。\footnote{这一相消是该一自由度集总模型的特性。积分力控对环境刚度的真实依赖（见下方洞察）须在含传动柔性的高阶模型上才显式浮现，参 Ferretti, Magnani \& Rocco（1995）对单关节四阶模型的分析。}

\begin{insight}[Raibert--Craig 的稳定性不是无条件的：积分力控对\emph{刚}环境更易失稳]
\label{ins:fw-rc-cond}
力控子空间的稳定性\emph{并非}无条件，且其与环境刚度 $K_e$ 的真实关系常被误记反向。含传动柔性的完整分析（Ferretti, Magnani \& Rocco 1995，单关节四阶模型）给出的结论是：\textbf{积分力控随接触刚度\emph{增大}而稳定裕度\emph{下降}}——即在\emph{软}环境上可用较大积分增益 $K_{fi}$，而碰\emph{刚}环境（钢铁、刚性夹具）时同样的 $K_{fi}$ 会激起振荡乃至失稳（这也是导纳式力控在刚性接触下变不稳的同一机理）。\textbf{这不是无条件稳定（不像阻抗控制有无源性保证），而是取决于环境刚度与力控增益的匹配}——这是混合方案的固有弱点：它需要对环境有一定先验（尤其在刚性接触侧）才能保证稳定。
\end{insight}

\paragraph{局限四表。}
\cref{tab:fw-rc-limits} 总结 Raibert--Craig 的四点根本局限。其中“需已知约束几何”尤为致命：选择矩阵 $\mathbf{S}_\sigma$ \emph{不是}工程师随意挑的设计参数，而是\textbf{由约束几何决定的物理量}（Mason 1981\cite{mason1981compliance}）。若约束帧方向估计偏差 $5^\circ$，本应控力的分量会“泄漏”到控位方向——在 $10\,\mathrm{N}$ 法向力下产生约 $10\sin5^\circ\approx0.87\,\mathrm{N}$ 的侧向干扰，对精密装配可能不可接受。

\begin{table}[H]\centering\small
\caption{Raibert--Craig 力位混合的四点局限}
\label{tab:fw-rc-limits}
\begin{tabular}{lll}
\toprule
局限 & 原因 & 后果 \\
\midrule
需已知约束几何 & $\mathbf{S}_\sigma$ 需精确约束帧方向 & 曲面/未知表面不适用 \\
不连续切换 & $s_i\in\{0,1\}$ 无中间态 & 接触/脱离时控制律跳变 \\
力位耦合 & 控力方向位移影响控位方向 & 交叉耦合致性能退化 \\
无阻抗行为 & 控力方向无弹簧/阻尼 & 接触瞬间力不连续 \\
\bottomrule
\end{tabular}
\end{table}

% ════════════════════════════════════════════════════════════════
\section{直接力控与抗饱和}
\label{sec:fw-direct}

当任务只关心\emph{精确跟踪一个力参考}（如恒力打磨），可用\emph{直接力控}：外环力 PI、内环位置/速度的级联结构。外环力误差 $\mathbf{e}_f=\mathbf{f}_d-\mathbf{f}$ 经 PI 生成位置参考增量给内环，内环负责快速跟踪。这套结构简单可靠，但有一个隐患——\textbf{积分漂移}。

\paragraph{积分漂移与三种抗饱和。}
当末端尚未接触（$\mathbf{f}\approx\mathbf{0}$）时，力误差 $\mathbf{e}_f\approx\mathbf{f}_d$ 持续非零，积分器 $\int\mathbf{e}_f$ 不断累积；一旦接触，巨大的积分量瞬间释放成位置参考跳变，导致冲击力尖峰。三种工业常用对策（\cref{tab:fw-antiwindup}）：\emph{积分限幅}（clamp）最简单但限值难定；\emph{条件积分}（gating）只在接触时积分，最干净；\emph{泄漏积分器} $\dot{I}=e-\alpha I$ 给积分一个一阶低通遗忘 $1/(s+\alpha)$，自动衰减历史，代价是稳态偏差 $I_{ss}=e_{ss}/\alpha$（需在“抗漂移”与“稳态精度”间权衡 $\alpha$）。

\begin{table}[H]\centering\small
\caption{直接力控的三种抗积分饱和策略}
\label{tab:fw-antiwindup}
\begin{tabular}{llll}
\toprule
策略 & 机制 & 优点 & 代价 \\
\midrule
积分限幅 & $I\!\leftarrow\!\mathrm{clip}(I,-I_{\max},I_{\max})$ & 实现极简 & 限值难定 \\
条件积分 & 仅接触时累积 & 无漂移、最干净 & 需可靠接触检测 \\
泄漏积分器 & $\dot I=e-\alpha I$（一阶低通） & 自动遗忘历史 & 稳态偏差 $e_{ss}/\alpha$ \\
\bottomrule
\end{tabular}
\end{table}

条件积分的门控逻辑可用一段语言无关的示意伪代码表达（这是全章仅有的两处示意代码之一）：
\begin{quote}\small\ttfamily
in\_contact $\leftarrow$ ($\|f\|>f_{\rm thresh}$)\\
if in\_contact: \quad I $\leftarrow$ I $+ e_f \cdot \Delta t$\quad // 仅接触时积分\\
else: \qquad\qquad\ \, I $\leftarrow$ 0\qquad\qquad\quad // 脱离即清零，杜绝漂移\\
$F_{\rm force} \leftarrow K_{fp}\, e_f + K_{fi}\, I$
\end{quote}

% ════════════════════════════════════════════════════════════════
\section{主动刚度控制（Salisbury）与各向异性柔顺}
\label{sec:fw-stiffness}

Salisbury（1980）的\emph{主动刚度控制}\cite{salisbury1980active} 是阻抗的简化前身：让末端表现为一个纯虚拟弹簧。控制律去掉了阻尼与惯量整形，只保留刚度项：
\begin{equation}
  \boldsymbol{\tau}=\mathbf{J}^\top\mathbf{K}_x\,(\mathbf{x}_d-\mathbf{x})+\mathbf{g}(\mathbf{q})
  =\mathbf{J}^\top\mathbf{K}_x\mathbf{e}_x+\mathbf{g}.
  \label{eq:fw-stiffness-law}
\end{equation}
在外力 $\mathbf{f}_{\rm ext}$ 下静态平衡时，$\mathbf{K}_x\tilde{\mathbf{x}}=\mathbf{f}_{\rm ext}$（$\tilde{\mathbf{x}}=\mathbf{x}-\mathbf{x}_d=-\mathbf{e}_x$），故位移
\begin{equation}
  \tilde{\mathbf{x}}=\mathbf{K}_x^{-1}\mathbf{f}_{\rm ext},
  \label{eq:fw-stiffness-spring}
\end{equation}
这正是胡克定律——末端偏离与外力成正比、与刚度成反比。

\paragraph{各向异性与构型相关的关节等效刚度。}
$\mathbf{K}_x$ 取对角即可在不同笛卡尔方向设\emph{各向异性}刚度（如装配中让插入方向软、横向硬）。把任务空间刚度投影回关节空间，得\emph{关节等效刚度}
\begin{equation}
  \mathbf{K}_q=\mathbf{J}^\top\mathbf{K}_x\mathbf{J}.
  \label{eq:fw-stiffness-joint}
\end{equation}
注意 $\mathbf{K}_q$ \emph{随构型变化}（$\mathbf{J}=\mathbf{J}(\mathbf{q})$）——同一笛卡尔刚度在不同位形下对应不同的关节刚度，这是任务空间柔顺设计相对关节空间的优势。

\begin{insight}[纯弹簧会弹跳——Hogan 加阻尼与惯量的动机]
\label{ins:fw-stiffness-bounce}
\cref{eq:fw-stiffness-law} 是\emph{无阻尼}的纯弹簧。取势能 $V=\tfrac12\tilde{\mathbf{x}}^\top\mathbf{K}_x\tilde{\mathbf{x}}$ 作 Lyapunov 函数可证系统\emph{Lyapunov 稳定}（有界），但 $\dot V$ 不严格负定，无法保证收敛——物理上表现为接触后持续\textbf{振荡甚至弹跳}（“纯弹簧手指”）。要让能量\emph{耗散}从而渐近稳定，必须补上阻尼项 $\mathbf{D}_d$（必要时再整形惯量 $\mathbf{M}_d$）——这正是下一步走向 Hogan 完整阻抗控制（\cref{sec:fw-impedance,sec:fw-passivity}）的直接动机。
\end{insight}

% ════════════════════════════════════════════════════════════════
\section{并行力/位控制（Chiaverini 1993）}
\label{sec:fw-parallel}

Raibert--Craig 用\emph{硬切换}的选择矩阵分管力与位置，需已知约束几何。Chiaverini \& Sciavicco（1993）\cite{chiaverini1993parallel} 提出另一思路：力 PI 与位 PD 在\emph{所有}方向上\emph{同时}作用，不用选择矩阵，靠\textbf{增益大小}做\emph{隐式}分解——力控方向把力增益设大、位控方向把位增益设大，二者在每个方向上“并行”竞争。其控制律为
\begin{equation}
  \boldsymbol{\tau}=\mathbf{J}^\top\!\Big[\underbrace{\mathbf{K}_p\mathbf{e}_x+\mathbf{K}_d\dot{\mathbf{e}}_x}_{\text{位 PD（全方向）}}+\underbrace{\mathbf{K}_{fp}\mathbf{e}_f+\mathbf{K}_{fi}\!\textstyle\int\mathbf{e}_f}_{\text{力 PI（全方向）}}\Big]+\mathbf{C}\dot{\mathbf{q}}+\mathbf{g}.
  \label{eq:fw-parallel-law}
\end{equation}

\paragraph{为何力控“赢”：积分作用主导。}
关键设计是让\emph{力环含积分、位环只含比例}。在受约束方向上，位 PD 试图维持位置参考，力 PI 试图维持力参考；二者的稳态博弈中，\textbf{力的积分项最终主导}——只要存在力误差，积分持续累积直到把位置“推”到使 $\mathbf{e}_f=\mathbf{0}$。把力环（含弹性环境 $f=K_e\tilde{x}$）的闭环力传函写出，由终值定理可证\emph{稳态力误差为零}：积分天然消除常值力误差，而位环的比例项在受约束方向上只贡献有限“抗力”、不阻止积分收敛。这就是“并行”而非“互斥”：力控方向的胜出靠积分作用而非显式屏蔽位控。

\begin{remark}[并行控制的接口量纲]
\label{rmk:fw-parallel-interface}
与 \cref{rmk:fw-interface} 同理，$\mathbf{K}_{fp},\mathbf{K}_{fi}$ 的物理意义取决于接口。当力 PI 输出位置偏置时，乘积 $\mathbf{K}_e\mathbf{K}_{fp}$ 具\emph{阻尼}量纲、$\mathbf{K}_e\mathbf{K}_{fi}$ 具\emph{刚度}量纲——并行律因此可解读为“在力反馈中隐式注入了一组与环境刚度耦合的弹簧--阻尼”。
\end{remark}

\cref{tab:fw-rc-vs-parallel} 对比两种方案。并行控制的最大优势是\emph{无需约束几何先验}，适合\textbf{未知曲面跟踪}（如打磨自由曲面）——它不必预先知道法向方向，靠力积分自动“贴合”表面。代价是它也是\emph{条件稳定}的（无无源性保证，稳定性依赖增益与 $K_e$ 匹配）。

\begin{table}[H]\centering\small
\caption{Raibert--Craig 力位混合 vs. Chiaverini 并行力/位}
\label{tab:fw-rc-vs-parallel}
\begin{tabular}{lll}
\toprule
维度 & Raibert--Craig（1981） & Chiaverini（1993） \\
\midrule
力位分解 & 显式（选择矩阵 $\mathbf{S}_\sigma$ 硬切换） & 隐式（增益大小，全方向并行） \\
约束几何 & 必须已知 & 无需先验 \\
稳态力误差 & 取决于力 PI & 零（终值定理，积分主导） \\
适用场景 & 已知平面/规则约束 & 未知曲面跟踪 \\
稳定性 & 条件稳定（依赖 $K_e$） & 条件稳定（依赖 $K_e$） \\
\bottomrule
\end{tabular}
\end{table}

% ════════════════════════════════════════════════════════════════
\section{阻抗控制完整推导}
\label{sec:fw-impedance}

至此的四种算法各有短板：力位混合需约束几何且硬切换、纯刚度会弹跳、并行控制无稳定性保证。Hogan 阻抗控制把它们统一并补足——\textbf{不分别控力或控位，而是整形末端的二阶力--运动关系}，自带阻尼故不弹跳、含无源性故对环境鲁棒。\cref{sec:ctrl-robot-impedance} 已给出目标阻抗 \cref{eq:ctrl-impedance} 与免力传感特例 \cref{eq:ctrl-impedance-law}；本节给出\emph{完整四步推导}与\emph{零空间}处理（比导论深一层），稳定性证明留到 \cref{sec:fw-passivity}。

\subsection{关节空间阻抗}
\label{sec:fw-impedance-joint}

最简单的版本是让\emph{每个关节}表现为虚拟弹簧--阻尼。期望行为 $\mathbf{M}_d^q\ddot{\mathbf{e}}+\mathbf{D}_d^q\dot{\mathbf{e}}+\mathbf{K}_d^q\mathbf{e}=\boldsymbol{\tau}_{\rm ext}$（$\mathbf{e}=\mathbf{q}-\mathbf{q}_d$）。从关节动力学 $\mathbf{M}\ddot{\mathbf{q}}+\mathbf{C}\dot{\mathbf{q}}+\mathbf{g}=\boldsymbol{\tau}+\boldsymbol{\tau}_{\rm ext}$ 出发，取 $\mathbf{M}_d^q\approx\mathbf{M}(\mathbf{q})$（不整形关节惯量）反解，得\emph{计算力矩法 + 阻抗项}：
\begin{equation}
  \boldsymbol{\tau}=\mathbf{M}\ddot{\mathbf{q}}_d+\mathbf{C}\dot{\mathbf{q}}+\mathbf{g}-\mathbf{K}_d^q(\mathbf{q}-\mathbf{q}_d)-\mathbf{D}_d^q(\dot{\mathbf{q}}-\dot{\mathbf{q}}_d).
  \label{eq:fw-imp-joint}
\end{equation}
工程中最常用的简化是\emph{重力补偿 + PD}：$\boldsymbol{\tau}=\mathbf{g}-\mathbf{K}_d^q(\mathbf{q}-\mathbf{q}_d)-\mathbf{D}_d^q\dot{\mathbf{q}}$，省掉惯量与科氏前馈。对低速接触（打磨、装配）已足够好；快速运动时非线性耦合会“污染”阻抗（某些方向显得比预期更硬或更软）。

\subsection{笛卡尔空间阻抗（含动力学补偿）}
\label{sec:fw-impedance-cart}

更有用的是直接在\emph{末端}整形阻抗。沿用 \cref{eq:ctrl-impedance} 的本书约定（$\mathbf{e}_x=\mathbf{x}_d-\mathbf{x}$、右端 $-\mathcal{F}_{\rm ext}$，其中 $\mathcal{F}_{\rm ext}$ 为环境施于末端的外力）。源文献常写 $\mathbf{M}_d\ddot{\tilde{\mathbf{x}}}+\mathbf{D}_d\dot{\tilde{\mathbf{x}}}+\mathbf{K}_d\tilde{\mathbf{x}}=\mathbf{f}_{\rm ext}$（$\tilde{\mathbf{x}}=\mathbf{x}-\mathbf{x}_d$），其中 $\mathbf{f}_{\rm ext}$ 与 $\mathcal{F}_{\rm ext}$ 是\emph{同一}外力（皆为环境施于末端者，$\mathbf{f}_{\rm ext}=\mathcal{F}_{\rm ext}$）；因 $\tilde{\mathbf{x}}=-\mathbf{e}_x$，对源式整体乘 $-1$ 即得本书式（符号差全部由 $\tilde{\mathbf{x}}=-\mathbf{e}_x$ 承担，右端不另变号）。取最常用的 $\mathbf{M}_d=\boldsymbol{\Lambda}$（\emph{不整形末端等效质量}，由此无需力传感），推导大幅简化：取最常用的 $\mathbf{M}_d=\boldsymbol{\Lambda}$（\emph{不整形末端等效质量}，由此无需力传感），推导大幅简化：

\begin{theorem}[笛卡尔阻抗控制律（$\mathbf{M}_d=\boldsymbol{\Lambda}$）]
\label{thm:fw-cart-impedance}
设操作空间动力学 $\boldsymbol{\Lambda}\ddot{\mathbf{x}}+\boldsymbol{\mu}+\mathbf{p}=\mathbf{F}$（\cref{eq:ctrl-osc}）。命令加速度取
\begin{equation}
  \ddot{\mathbf{x}}_{\rm cmd}=\ddot{\mathbf{x}}_d+\boldsymbol{\Lambda}^{-1}\!\big(\mathbf{D}_d\dot{\mathbf{e}}_x+\mathbf{K}_d^x\mathbf{e}_x\big),
  \label{eq:fw-cart-xcmd}
\end{equation}
经操作空间动力学实现为任务力 $\mathbf{F}_{\rm cmd}=\boldsymbol{\Lambda}\ddot{\mathbf{x}}_{\rm cmd}+\boldsymbol{\mu}+\mathbf{p}=\boldsymbol{\Lambda}\ddot{\mathbf{x}}_d+\mathbf{D}_d\dot{\mathbf{e}}_x+\mathbf{K}_d^x\mathbf{e}_x+\boldsymbol{\mu}+\mathbf{p}$，再经对偶映射 $\boldsymbol{\tau}=\mathbf{J}^\top\mathbf{F}_{\rm cmd}$（含零空间项），即令闭环呈现目标阻抗 $\boldsymbol{\Lambda}\ddot{\mathbf{e}}_x+\mathbf{D}_d\dot{\mathbf{e}}_x+\mathbf{K}_d^x\mathbf{e}_x=-\mathcal{F}_{\rm ext}$。
\end{theorem}

\begin{proof}[骨架]
真实操作空间动力学受外力为 $\boldsymbol{\Lambda}\ddot{\mathbf{x}}+\boldsymbol{\mu}+\mathbf{p}=\mathbf{F}_{\rm cmd}+\mathcal{F}_{\rm ext}$（环境外力 $\mathcal{F}_{\rm ext}$ 与作动力 $\mathbf{F}_{\rm cmd}$ 同侧作用于末端，故在任务方程中同号叠加）。代入 $\mathbf{F}_{\rm cmd}$ 并消去 $\boldsymbol{\mu}+\mathbf{p}$：$\boldsymbol{\Lambda}\ddot{\mathbf{x}}=\boldsymbol{\Lambda}\ddot{\mathbf{x}}_d+\mathbf{D}_d\dot{\mathbf{e}}_x+\mathbf{K}_d^x\mathbf{e}_x+\mathcal{F}_{\rm ext}$。两边减 $\boldsymbol{\Lambda}\ddot{\mathbf{x}}_d$ 并用 $\ddot{\mathbf{e}}_x=\ddot{\mathbf{x}}_d-\ddot{\mathbf{x}}$（故 $\ddot{\mathbf{x}}-\ddot{\mathbf{x}}_d=-\ddot{\mathbf{e}}_x$），整理得 $\boldsymbol{\Lambda}\ddot{\mathbf{e}}_x+\mathbf{D}_d\dot{\mathbf{e}}_x+\mathbf{K}_d^x\mathbf{e}_x=-\mathcal{F}_{\rm ext}$，即目标阻抗。$\boldsymbol{\Lambda}$ 在 $\mathbf{F}_{\rm cmd}$ 与右端\emph{同时出现}是 $\mathbf{M}_d=\boldsymbol{\Lambda}$ 的红利：外力项系数化为 $\pm\mathbf{I}$（移到误差侧后为 $-\mathbf{I}$），无需测量 $\mathcal{F}_{\rm ext}$。
\end{proof}

\paragraph{静态参考下的实用形式与零空间。}
在静态参考（$\dot{\mathbf{x}}_d=\ddot{\mathbf{x}}_d=\mathbf{0}$）下，\cref{thm:fw-cart-impedance} 的力矩落为冗余臂上的实用律：
\begin{equation}
  \boxed{\;\boldsymbol{\tau}=\mathbf{J}^\top\!\big(\mathbf{K}_d^x\mathbf{e}_x-\mathbf{D}_d\,\mathbf{J}\dot{\mathbf{q}}\big)+\big(\mathbf{I}-\mathbf{J}^\top\bar{\mathbf{J}}^\top\big)\boldsymbol{\tau}_{\rm null}+\mathbf{C}\dot{\mathbf{q}}+\mathbf{g}.\;}
  \label{eq:fw-cart-law}
\end{equation}
阻尼项写成 $-\mathbf{D}_d\mathbf{J}\dot{\mathbf{q}}$ 而非 $-\mathbf{D}_d\dot{\tilde{\mathbf{x}}}$，是为只用可测的 $\dot{\mathbf{q}}$、免去对 $\mathbf{x}_d$ 做带噪数值微分。

\begin{remark}[零空间须用动力学一致伪逆，否则泄漏到末端]
\label{rmk:fw-nullspace}
\cref{eq:fw-cart-law} 的零空间投影 $(\mathbf{I}-\mathbf{J}^\top\bar{\mathbf{J}}^\top)$ 用的是\emph{动力学一致}逆 $\bar{\mathbf{J}}=\mathbf{M}^{-1}\mathbf{J}^\top\boldsymbol{\Lambda}$——它是\emph{唯一}使任意零空间力矩对末端\emph{零加速度}的广义逆（Khatib 1987\cite{khatib1987operational}，详见 O 章 \cref{ch:operational-space} 与 \cref{sec:ctrl-robot-osc}）。许多工业实现（如 libfranka 的笛卡尔阻抗示例）图省事用 Moore--Penrose 逆 $\mathbf{J}^+$ 构造 $(\mathbf{I}-\mathbf{J}^\top\mathbf{J}^{+\top})$，这是纯运动学投影、\emph{非}动力学一致——后果是零空间的关节中心化/避奇异力矩会“泄漏”出一个末端扰动加速度，破坏阻抗的纯净。这是 \cref{rmk:force-wbc-notation} 强调“两种伪逆严格区分”的工程缘由。
\end{remark}

\begin{insight}[阻抗的真正读法：弹簧--阻尼，刚度是柔顺旋钮]
\label{ins:fw-impedance-read}
\cref{eq:fw-cart-law} 与 \cref{sec:ctrl-robot-osc} 的解耦律本是同一物。差别只在\emph{读法}：解耦律读作“前馈 + 任务空间 PD + 自然惯量”，阻抗读作“质量--弹簧--阻尼”。把 $\mathbf{K}_d^x$ 调小末端就“软”（顺从外力、安全交互），调大就“硬”（精确定位、抗扰）——同一控制器靠一个参数从位置控制\emph{连续}过渡到柔顺控制。这把前四节“非此即彼”的力/位之争，化为一根可连续旋拧的旋钮。
\end{insight}

% ════════════════════════════════════════════════════════════════
\section{阻抗控制的稳定性：无源性与 LaSalle}
\label{sec:fw-passivity}

阻抗控制之所以成为力控的“默认选择”，根在它的稳定性\emph{最完善}：不仅自身渐近稳定，还经无源性保证与\emph{任何}被动环境耦合后仍稳定。本节给出从存储函数到渐近稳定的完整链条（与 \cref{ch:lyapunov-stability} 的 Lyapunov/LaSalle 工具互引）。

\begin{theorem}[笛卡尔阻抗的无源性与渐近稳定]
\label{thm:fw-passivity}
对目标阻抗 $\mathbf{M}_d\ddot{\tilde{\mathbf{x}}}+\mathbf{D}_d\dot{\tilde{\mathbf{x}}}+\mathbf{K}_d\tilde{\mathbf{x}}=\mathbf{f}_{\rm ext}$（$\mathbf{M}_d,\mathbf{D}_d,\mathbf{K}_d\succ0$），取存储函数
\begin{equation}
  V=\tfrac12\dot{\tilde{\mathbf{x}}}^\top\mathbf{M}_d\dot{\tilde{\mathbf{x}}}+\tfrac12\tilde{\mathbf{x}}^\top\mathbf{K}_d\tilde{\mathbf{x}}\;(=\text{动能}+\text{势能}),
  \label{eq:fw-storage}
\end{equation}
则 $\dot V=\dot{\tilde{\mathbf{x}}}^\top\mathbf{f}_{\rm ext}-\dot{\tilde{\mathbf{x}}}^\top\mathbf{D}_d\dot{\tilde{\mathbf{x}}}$（严格输出无源）。无外力（$\mathbf{f}_{\rm ext}=\mathbf{0}$）时原点\emph{全局渐近稳定}；有常值外力时稳态偏移 $\tilde{\mathbf{x}}_{ss}=\mathbf{K}_d^{-1}\mathbf{f}_{{\rm ext},ss}$（弹簧行为）。
\end{theorem}

\begin{proof}[完整链条]
\textbf{(A) 构造与 (B) 正定径向无界}：$\mathbf{M}_d,\mathbf{K}_d\succ0$ 保证 $V>0\,(\forall(\tilde{\mathbf{x}},\dot{\tilde{\mathbf{x}}})\neq\mathbf{0})$、$V=0\!\iff\!(\tilde{\mathbf{x}},\dot{\tilde{\mathbf{x}}})=\mathbf{0}$，且 $V\!\to\!\infty$ 当 $\|(\tilde{\mathbf{x}},\dot{\tilde{\mathbf{x}}})\|\!\to\!\infty$（径向无界 $\Rightarrow$ 轨迹有界）。
\textbf{(C) 求 $\dot V$}：$\dot V=\dot{\tilde{\mathbf{x}}}^\top\mathbf{M}_d\ddot{\tilde{\mathbf{x}}}+\tilde{\mathbf{x}}^\top\mathbf{K}_d\dot{\tilde{\mathbf{x}}}$，代入 $\mathbf{M}_d\ddot{\tilde{\mathbf{x}}}=\mathbf{f}_{\rm ext}-\mathbf{D}_d\dot{\tilde{\mathbf{x}}}-\mathbf{K}_d\tilde{\mathbf{x}}$，得 $\dot V=\dot{\tilde{\mathbf{x}}}^\top\mathbf{f}_{\rm ext}-\dot{\tilde{\mathbf{x}}}^\top\mathbf{D}_d\dot{\tilde{\mathbf{x}}}-\dot{\tilde{\mathbf{x}}}^\top\mathbf{K}_d\tilde{\mathbf{x}}+\tilde{\mathbf{x}}^\top\mathbf{K}_d\dot{\tilde{\mathbf{x}}}$；后两项因 $\mathbf{K}_d$ 对称（$\mathbf{a}^\top\mathbf{K}_d\mathbf{b}=\mathbf{b}^\top\mathbf{K}_d\mathbf{a}$）\emph{相消}。
\textbf{(D) 符号}：$\mathbf{f}_{\rm ext}=\mathbf{0}$ 时 $\dot V=-\dot{\tilde{\mathbf{x}}}^\top\mathbf{D}_d\dot{\tilde{\mathbf{x}}}\le0$，且 $\dot V=0\!\iff\!\dot{\tilde{\mathbf{x}}}=\mathbf{0}$（因 $\mathbf{D}_d\succ0$）——\emph{半}负定。
\textbf{(E) LaSalle}：系统收敛到 $\{\dot V=0\}=\{\dot{\tilde{\mathbf{x}}}=\mathbf{0}\}$ 中最大不变集；在其上 $\ddot{\tilde{\mathbf{x}}}=\mathbf{0}$，代回动力学得 $\mathbf{K}_d\tilde{\mathbf{x}}=\mathbf{0}\Rightarrow\tilde{\mathbf{x}}=\mathbf{0}$。故最大不变集 $=\{\mathbf{0}\}$，原点全局渐近稳定。
\end{proof}

\begin{insight}[$\dot V\le0$ 只给 Lyapunov 稳定，须 LaSalle 才到渐近稳定]
\label{ins:fw-lasalle}
许多推导在 \emph{Step D}（得 $\dot V\le0$）后径直宣称“渐近稳定”——这是\textbf{常见扣分点}。$\dot V\le0$ 只保证 Lyapunov 稳定（有界），因为在 $\dot{\tilde{\mathbf{x}}}=\mathbf{0},\tilde{\mathbf{x}}\neq\mathbf{0}$ 处 $\dot V=0$ 而系统未必停在那里。\emph{必须}用 LaSalle 不变集定理（\emph{Step E}）排除“振荡不衰减”的可能，才能升级到渐近稳定。同理，若把 $\mathbf{D}_d$ 设为\emph{半}正定（某方向阻尼为零，如忘了给某旋转方向设阻尼），那个方向就仅 Lyapunov 稳定而持续振荡——\textbf{每个自由度方向（含旋转）都必须有正阻尼}。
\end{insight}

\paragraph{Colgate--Hogan 耦合稳定定理：对环境刚度的根本鲁棒性。}
\cref{thm:fw-passivity} 是机器人\emph{孤立}时的结论。与环境接触时，真正要保证的是\emph{耦合系统}稳定。这正是无源性的威力所在：

\begin{theorem}[Colgate--Hogan 1988 耦合稳定]
\label{thm:fw-colgate}
若控制器（机械手）是输出无源的、且环境是无源的（被动，$\mathbf{K}_e\succ0$），则\emph{耦合系统稳定}\cite{colgate1988robust}。
\end{theorem}

其深刻含义是：阻抗控制器\emph{不需要“知道”环境刚度} $K_e$ 就能保持稳定——它靠自身无源性“自动”适配任何被动环境。这与前几节算法形成鲜明对比：力 PI（\cref{eq:fw-rc-rh}）、混合控制、并行控制的稳定条件都\emph{显式含} $K_e$，故对环境敏感（\cref{tab:fw-stability}）。阻抗对从泡沫（$K_e\!\sim\!10^2$）到钢铁（$K_e\!>\!10^6$）的全谱环境都稳定，这正是它成为交互控制基石的理论根由。

% ════════════════════════════════════════════════════════════════
\section{柔性关节阻抗与自适应力控}
\label{sec:fw-flex-adaptive}

前述推导默认关节\emph{刚性}。但 Franka Panda、KUKA iiwa 等协作臂的关节含\emph{弹性}（谐波减速器 + 力矩传感器弹性），不能当刚性处理；且真实动力学参数（负载质量、摩擦）常\emph{未知}。本节补两块：柔性关节阻抗（Ott 2008）与自适应力控（Slotine--Li 1987）。

\subsection{柔性关节阻抗（Ott 2008）}
\label{sec:fw-flex}

把柔性关节建模为\emph{双质量}系统：电机侧惯量 $B$、连杆侧惯量 $\mathbf{M}_l$、二者由刚度 $K_s$ 的关节弹簧连接，弹簧力矩 $\tau_s=K_s(q_m-q_l)$。控制策略是\emph{内外双环}：内环（电机侧）整形为高带宽位置环，外环（连杆侧）实现期望笛卡尔阻抗。Ott 的核心结论是一个\emph{物理硬上界}：

\begin{theorem}[等效笛卡尔刚度上界]
\label{thm:fw-ott}
柔性关节机器人可实现的等效笛卡尔刚度满足
\begin{equation}
  K_{d,eq}=\frac{K_\theta K_s}{K_\theta+K_s}\le K_s,
  \label{eq:fw-ott-bound}
\end{equation}
其中 $K_\theta$ 为外环设定刚度、$K_s$ 为关节弹簧刚度\cite{ott2008cartesian,ott2008unified}。即\textbf{用户设定的刚度不能超过关节弹簧刚度}——这是物理限制而非软件限制。
\end{theorem}

\cref{eq:fw-ott-bound} 是两弹簧串联公式：$K_\theta$ 与 $K_s$ 串联，等效刚度由\emph{较软}者主导，恒 $\le K_s$。若忽视它、强行设 $K_d>K_s$（让连杆侧“硬”过关节弹簧），物理上不可能，结果是关节弹簧振荡、力矩饱和、安全停机。这就是 Franka/iiwa 控制器为何天然限幅刚度的理论边界。\cref{eq:fw-ott-bound} 也解释了驱动选型的取舍（\cref{tab:fw-sea-qdd}）：\emph{串联弹性驱动}（SEA，如 Franka/iiwa）$K_s$ 有限，故可实现刚度与力控带宽受限，但弹簧天然吸收冲击、利于碰撞安全与力矩感知；\emph{准直驱}（QDD，如 MIT Cheetah）$K_s$ 极大（近刚性），可实现高刚度与高带宽，代价是碰撞保护弱。

\begin{table}[H]\centering\small
\caption{SEA 与 QDD 的力控取舍（\cref{eq:fw-ott-bound} 的工程含义）}
\label{tab:fw-sea-qdd}
\begin{tabular}{llll}
\toprule
驱动类型 & 关节弹簧 $K_s$ & 可实现笛卡尔刚度 & 力控带宽 \\
\midrule
SEA（Franka/iiwa） & 有限（中等） & 受 $K_s$ 上界限制 & 受限（弹性反共振） \\
QDD（MIT Cheetah） & 极大（近刚性） & 高 & 高 \\
\bottomrule
\end{tabular}
\end{table}

\subsection{自适应力控（Slotine--Li 1987）}
\label{sec:fw-adaptive}

当动力学参数 $\boldsymbol{\theta}$（质量、质心、摩擦）\emph{未知}时，用在线辨识替代固定参数。机器人动力学的关键性质是\textbf{线性参数化}：无论参数取值如何，$\mathbf{M}\ddot{\mathbf{q}}_r+\mathbf{C}\dot{\mathbf{q}}_r+\mathbf{g}=\mathbf{Y}(\mathbf{q},\dot{\mathbf{q}},\dot{\mathbf{q}}_r,\ddot{\mathbf{q}}_r)\boldsymbol{\theta}$，回归矩阵 $\mathbf{Y}$ 与参数无关（重力的最简特例 $\mathbf{g}(\mathbf{q})=\mathbf{Y}_g(\mathbf{q})\boldsymbol{\theta}_g$）。Slotine--Li\cite{slotine1987adaptive} 框架引入\emph{参考速度}与\emph{滑模变量}：
\begin{equation}
  \dot{\mathbf{q}}_r=\dot{\mathbf{q}}_d-\alpha(\mathbf{q}-\mathbf{q}_d),\qquad
  \mathbf{s}=\dot{\mathbf{q}}-\dot{\mathbf{q}}_r=\dot{\mathbf{e}}+\alpha\mathbf{e}\quad(\alpha>0),
  \label{eq:fw-sliding}
\end{equation}
当 $\mathbf{s}\to\mathbf{0}$ 时 $\dot{\mathbf{e}}=-\alpha\mathbf{e}$ 使误差指数收敛。控制律 + 更新律为
\begin{equation}
  \boldsymbol{\tau}=\mathbf{Y}\hat{\boldsymbol{\theta}}+\mathbf{J}^\top\!\big(\mathbf{K}_d^x\mathbf{e}_x-\mathbf{D}_d\dot{\tilde{\mathbf{x}}}\big),\qquad
  \dot{\hat{\boldsymbol{\theta}}}=-\boldsymbol{\Gamma}\mathbf{Y}^\top\mathbf{s},
  \label{eq:fw-adaptive-law}
\end{equation}
$\boldsymbol{\Gamma}\succ0$ 为学习率矩阵。

\begin{theorem}[自适应力控的跟踪收敛]
\label{thm:fw-adaptive}
取 $V=\tfrac12\mathbf{s}^\top\mathbf{M}\mathbf{s}+\tfrac12\tilde{\mathbf{x}}^\top\mathbf{K}_d\tilde{\mathbf{x}}+\tfrac12\tilde{\boldsymbol{\theta}}^\top\boldsymbol{\Gamma}^{-1}\tilde{\boldsymbol{\theta}}$（$\tilde{\boldsymbol{\theta}}=\hat{\boldsymbol{\theta}}-\boldsymbol{\theta}$），则 $\dot V\le-\mathbf{s}^\top\mathbf{D}_d\mathbf{s}\le0$，由 Barbalat 引理 $\mathbf{s}\to\mathbf{0}$（跟踪误差渐近趋零）；若沿轨迹满足持续激励（PE）条件，则 $\tilde{\boldsymbol{\theta}}\to\mathbf{0}$（参数收敛真值）。
\end{theorem}

\begin{proof}[要点]
对 $V$ 求导，$\mathbf{s}^\top\mathbf{M}\dot{\mathbf{s}}+\tfrac12\mathbf{s}^\top\dot{\mathbf{M}}\mathbf{s}$ 中利用 $\dot{\mathbf{M}}-2\mathbf{C}$ 的\emph{反对称}性 $\mathbf{s}^\top(\dot{\mathbf{M}}-2\mathbf{C})\mathbf{s}=0$（该性质主放动力学部 \cref{ch:geometric-mechanics}，本章用结论）消去科氏项；代入控制律得误差动力学 $\mathbf{M}\dot{\mathbf{s}}+\mathbf{C}\mathbf{s}=\mathbf{Y}\tilde{\boldsymbol{\theta}}-(\text{阻抗反馈})$，其中 $\mathbf{s}^\top\mathbf{Y}\tilde{\boldsymbol{\theta}}$ 恰被更新律 $\dot{\hat{\boldsymbol{\theta}}}=-\boldsymbol{\Gamma}\mathbf{Y}^\top\mathbf{s}$ 贡献的 $\tilde{\boldsymbol{\theta}}^\top\boldsymbol{\Gamma}^{-1}\dot{\hat{\boldsymbol{\theta}}}=-\tilde{\boldsymbol{\theta}}^\top\mathbf{Y}^\top\mathbf{s}$ \emph{精确抵消}，余下 $-\mathbf{s}^\top\mathbf{D}_d\mathbf{s}$。系统非自治，用 Barbalat（非 LaSalle）下结论。
\end{proof}

\begin{insight}[无 PE 仍跟踪收敛，参数未必收敛——“等效参数”现象]
\label{ins:fw-pe}
自适应控制的精妙在于：即便参数\emph{永不}收敛到真值（PE 不满足），跟踪误差 $\mathbf{e}$ 仍渐近趋零。原因藏在 \cref{thm:fw-adaptive} 的存储函数里——第三项（参数误差能量）只保证\emph{不增长}、不保证收敛，而前两项（跟踪误差）确实收敛。\textbf{控制器可能用“错误”的参数实现正确的行为}：它找到的不是真实参数，而是一组在当前轨迹上恰好产生正确力矩的\emph{等效参数}。工程上须配\emph{参数投影}（如约束 $\hat m\in[0.01,50]\,\mathrm{kg}$）防参数漂到非物理负值。自适应与鲁棒/DOB/RL 残差的取舍详见辨识/鲁棒部 \cref{ch:robust-hinf}。
\end{insight}

% ════════════════════════════════════════════════════════════════
\section{稳定性方法统一与力控带宽}
\label{sec:fw-stability-bw}

\paragraph{五算法稳定性谱。}
\cref{tab:fw-stability} 把本章五大算法的稳定性\emph{类型}并列。一个清晰的梯度是：基于\emph{选择/增益分管}的算法（Raibert--Craig、并行）是\emph{条件稳定}（依赖 $K_e$）；纯刚度仅 Lyapunov 稳定（无阻尼）；唯有阻抗控制是\emph{渐近稳定 + 无源性保证}（经 Colgate--Hogan 对全谱环境稳定）。这解释了为何阻抗是“默认选择”。一切力控稳定性分析都可纳入统一的 Lyapunov 三步流程：\emph{(1)} 构造正定径向无界的存储函数（能量）；\emph{(2)} 求 $\dot V$ 并代入闭环；\emph{(3)} 由 $\dot V\le0$ 定 Lyapunov 稳定，必要时用 LaSalle/Barbalat 升级到渐近稳定。

\begin{table}[H]\centering\small
\caption{五大经典力控算法的稳定性类型对照}
\label{tab:fw-stability}
\begin{tabular}{llll}
\toprule
算法 & 分析方法 & 稳定性类型 & 对 $K_e$ 敏感 \\
\midrule
Raibert--Craig & 分通道 / Routh--Hurwitz & 条件稳定（需正确 $\mathbf{S}_\sigma$） & 是 \\
直接力控 PI & 闭环极点 & BIBO（适当增益） & 是 \\
主动刚度 & Lyapunov（弹性势能） & 仅 Lyapunov（无渐近） & 中 \\
并行力/位 & 闭环传函 / 终值定理 & 条件稳定 & 是 \\
阻抗控制 & 无源性 + LaSalle & \textbf{渐近 + 无源保证} & \textbf{否} \\
\bottomrule
\end{tabular}
\end{table}

\paragraph{离散化稳定条件。}
连续稳定不蕴含离散稳定：采样与零阶保持引入相位滞后。对 1-DOF 接触系统 $M_d\ddot{\tilde x}+D_d\dot{\tilde x}+(K_d+K_e)\tilde x=0$ 用前向 Euler，稳定要求
\begin{equation}
  \Delta t<\frac{2D_d}{K_d+K_e}.
  \label{eq:fw-discrete}
\end{equation}
$K_e$ 越大（刚性环境）许可步长越小——这是“高刚度环境需高采样率”的定量根由（经验上取 $f_{\rm ctrl}>5\sqrt{(K_d+K_e)/M_d}/2\pi$ 留 $5$ 倍过采样裕度）。

\paragraph{力控带宽：木桶效应与根本矛盾。}
力控环的\emph{有效带宽}由最慢环节决定（木桶效应，\cref{fig:fw-bandwidth}）：力传感器滤波、控制器计算延迟、关节柔性反共振、位置内环带宽、采样 Nyquist 限制层层串联，有效带宽 $=$ 诸环节之\emph{最小}。

\begin{figure}[H]
\centering
\begin{tikzpicture}[font=\footnotesize, node distance=2.2mm,
  band/.style={draw, rounded corners, align=center, minimum height=7mm, fill=main!8, draw=main!55!black, inner sep=3pt},
  ar/.style={->, thick, gray!55!black}]
  \node[band] (a) {力传感器滤波\\（最慢档之一）};
  \node[band, right=of a] (b) {控制计算延迟};
  \node[band, right=of b] (c) {关节柔性反共振};
  \node[band, right=of c] (d) {位置内环带宽\\（导纳的瓶颈）};
  \node[band, right=of d] (e) {Nyquist 限制};
  \node[band, right=of e, fill=second!10, draw=second!55!black] (f) {电流环\\（最快，非瓶颈）};
  \draw[ar] (a)--(b); \draw[ar] (b)--(c); \draw[ar] (c)--(d); \draw[ar] (d)--(e); \draw[ar] (e)--(f);
  \node[below=4mm of c.south, anchor=north, font=\footnotesize\itshape, xshift=8mm]
    {有效带宽 $=\min(\text{所有环节})$，由最短板决定。};
\end{tikzpicture}
\caption{力控带宽瓶颈链（木桶效应）。有效带宽由串联诸环节中最慢者封顶；电流环虽快却非瓶颈。}
\label{fig:fw-bandwidth}
\end{figure}

\begin{insight}[导纳带宽由内环位置环决定，而非导纳参数；力控与位控的根本矛盾]
\label{ins:fw-bw}
两个易错的本质洞察。\textbf{其一}：导纳控制的带宽\emph{不是}由导纳参数 $\mathbf{M}_d,\mathbf{D}_d,\mathbf{K}_d$ 决定的，\emph{而是}由\textbf{内环位置控制器带宽}决定——即便设计了高带宽导纳滤波器，若位置内环只有中等带宽，实际力控带宽仍被其封顶。这是导纳控制在刚性工业臂（UR、Fanuc）上性能受限的根本。\textbf{其二}：力控要\emph{柔顺}（低刚度、对力敏感），位控要\emph{刚性}（高刚度、抗扰），二者在同一方向\emph{矛盾}：
\begin{equation}
  \text{位控带宽}\propto\sqrt{K_{\rm total}/M},\qquad \text{力控灵敏度}\propto 1/K_{\rm total}.
  \label{eq:fw-bw-tradeoff}
\end{equation}
增大刚度提位控带宽却降力控灵敏度。力控设计的真正难点因此不在“选什么算法”，而在\textbf{随任务阶段动态调权}（自由空间硬、接触瞬间软）——这正是变阻抗与学习型力控要解决的核心。\emph{工程数字（微秒级计算、几十至上百赫兹带宽）随硬件/求解器浮动，此处只取定性量级。}
\end{insight}

至此，固定基座力控的完整谱系（\cref{sec:fw-raibert,,sec:fw-passivity,,sec:fw-stability-bw}）已铺成。它们共有一个隐含前提：\emph{基座固定、所有关节有驱动}。一旦机器人“站到冰面上”——基座浮动、欠驱动、多接触、多任务——这套框架就需要升维。下半章转入浮动基座的全身控制。

% ════════════════════════════════════════════════════════════════
\section{为何机械臂也需要 WBC}
\label{sec:fw-why-wbc}

固定基座 $7$-DoF 臂上，$\boldsymbol{\tau}=\mathbf{J}^\top(\mathbf{K}_d\mathbf{e}_x+\mathbf{D}_d\dot{\mathbf{e}}_x)+\mathbf{g}$ 已能做柔顺力/位控。但移动操作（底盘 + 臂）、人形上肢（浮动躯干 + 臂）、四足 + 臂这三类场景打破了上述前提，带来三大根本困难。

\paragraph{困难一：欠驱动的浮动基座。}
浮动基座有 $6$ 个自由度（$3$ 平移 $+3$ 旋转）却无直接执行器。浮基动力学（已在 \cref{sec:ctrl-robot-wbc} 的 \cref{eq:ctrl-floatbase} 与 \cref{ch:quad_wbc} 推导，此处只回顾结论）
\begin{equation}
  \mathbf{M}(\mathbf{q})\dot{\boldsymbol{\nu}}+\mathbf{C}\boldsymbol{\nu}+\boldsymbol{\tau}_g=\mathbf{S}^\top\boldsymbol{\tau}+\mathbf{J}_c^\top\boldsymbol{\lambda},\qquad
  \mathbf{S}^\top=\begin{bmatrix}\mathbf{0}_{6\times n_j}\\ \mathbf{I}_{n_j}\end{bmatrix},
  \label{eq:fw-floatbase}
\end{equation}
其\emph{前 $6$ 行}（基座方程）右端只有接触力 $\mathbf{J}_c^\top\boldsymbol{\lambda}$、没有 $\boldsymbol{\tau}$——\textbf{你无法直接给浮动基座施加力矩，只能经地面接触力间接控制}（位形 $\mathbf{q}=(\mathbf{q}_b,\mathbf{q}_j)$ 且 $n_q\neq n_v$，分块 $\mathbf{M}_{bb}/\mathbf{M}_{bj}$ 见 \cref{ch:reduced-models} 与 \cref{ch:quad_wbc}）。形象地说：固定基座臂像站在地板上推墙（脚钉死、手力直接来自手臂）；浮基机器人像\emph{站在冰面上推墙}（必须靠脚与地面的摩擦间接产生推力，否则自己被推走）——WBC 就是解“冰面推墙”的框架。

\paragraph{困难二：多任务冲突。}
固定基座阻抗只跟踪一个末端位姿；浮基机器人却要\emph{同时}协调一串带优先级的任务（动力学可行性 $>$ 摩擦锥 $>$ 质心平衡 $>$ 末端操作 $>$ 躯干姿态 $>$ 关节正则化），它们常相互冲突（末端想伸远但会让质心出支撑域）。固定基座阻抗没有处理多任务优先级的机制。

\paragraph{困难三：接触力须被优化。}
固定基座力控不关心接触力（基座螺栓承受所有反力）；浮基则不同——地面接触力是控制输出的一部分，必须满足摩擦锥（否则脚滑）与法向力非负（否则脚离地变自由落体）。\textbf{接触力不是“可选优化目标”，而是“必须满足的安全约束”}。摩擦锥的金字塔线性化（内/外逼近）已在 \cref{sec:ctrl-robot-wbc} 与接触力学部 \cref{ch:contact-mechanics}、\cref{ch:quad_wbc} 讲透，本章引用复用：轴对齐 $4$ 面金字塔 $|f_x|\le\mu f_z,|f_y|\le\mu f_z$ 是圆锥的\emph{外}逼近（box，非保守、不能当“保证不滑”的硬约束），内逼近用 $|f_x|+|f_y|\le\mu f_z$ 或等价 $\mu_{\rm eff}=\mu/\sqrt2$。

这三大困难指向同一答案——把力控变成一个\emph{带约束的多任务优化}。

% ════════════════════════════════════════════════════════════════
\section{面接触与 12 维顶点力（TSID Contact6d）}
\label{sec:fw-contact6d}

\cref{ch:quad_wbc} 处理腿足时把每个\emph{点}接触建模为 $3$D 力。但人形脚掌、四足的面接触脚是\emph{有面积}的——它传递一个完整 $6$D wrench $\mathbf{w}=[\mathbf{f};\boldsymbol{\tau}]$，且须约束\emph{压力中心}（CoP/ZMP）落在脚底矩形内。本节补 \cref{ch:quad_wbc} 缺的\emph{面接触}理论：TSID 的 $12$ 维顶点力参数化。

\paragraph{为何 6D wrench 不宜做决策变量。}
若直接以 $6$D wrench 为决策变量，摩擦锥仍是 SOCP；更糟的是 ZMP 约束 $-l_x\le-\tau_y/f_z\le l_x,\ -l_y\le\tau_x/f_z\le l_y$ 含 $\tau/f$ 的\emph{除法}，是\emph{非线性}约束，QP 无法直接处理。

\paragraph{TSID 的参数化。}
TSID 改用矩形脚底 $4$ 个顶点的力来\emph{参数化} wrench。顶点 $\mathbf{p}_i\,(i=1,\dots,4)$ 各承一个 $3$D 力 $\mathbf{f}_i$，决策变量 $\mathbf{f}_{12}=[\mathbf{f}_1;\mathbf{f}_2;\mathbf{f}_3;\mathbf{f}_4]\in\mathbb{R}^{12}$，经力生成矩阵映射为 wrench：
\begin{equation}
  \mathbf{w}=\mathbf{G}\,\mathbf{f}_{12},\qquad
  \mathbf{G}=\begin{bmatrix}\mathbf{I}_3 & \mathbf{I}_3 & \mathbf{I}_3 & \mathbf{I}_3\\[2pt] [\mathbf{p}_1]_\times & [\mathbf{p}_2]_\times & [\mathbf{p}_3]_\times & [\mathbf{p}_4]_\times\end{bmatrix}\in\mathbb{R}^{6\times12},
  \label{eq:fw-G}
\end{equation}
上块叠加合力、下块由 $\boldsymbol{\tau}=\sum_i\mathbf{p}_i\times\mathbf{f}_i$ 给合力矩（$[\cdot]_\times$ 为叉积反对称矩阵，见 \cref{ch:rigid_body}），$\mathrm{rank}(\mathbf{G})=6$。这一参数化的\emph{核心红利}是把非线性的 ZMP 约束化为线性的法向力非负约束：

\begin{theorem}[ZMP 凸包隐式满足]
\label{thm:fw-zmp}
若 $4$ 个顶点的法向力均非负（$f_{i,z}\ge0$），则 ZMP 自动落在 $4$ 顶点凸包（即脚底矩形）内，无需显式 ZMP 约束。
\end{theorem}

\begin{proof}
ZMP 是顶点按法向力加权的重心：
\begin{equation}
  \mathrm{ZMP}_x=\frac{\sum_i p_{i,x}f_{i,z}}{\sum_i f_{i,z}},\qquad
  \mathrm{ZMP}_y=\frac{\sum_i p_{i,y}f_{i,z}}{\sum_i f_{i,z}}.
  \label{eq:fw-zmp-cvx}
\end{equation}
当 $f_{i,z}\ge0$ 且 $\sum_i f_{i,z}>0$ 时，权重 $w_i=f_{i,z}/\sum_j f_{j,z}\ge0$、$\sum_i w_i=1$，故 $\mathrm{ZMP}=\sum_i w_i\mathbf{p}_i$ 是 $\{\mathbf{p}_i\}$ 的\emph{凸组合}，必落于其凸包。
\end{proof}

于是每个顶点只需独立的线性摩擦锥（$|f_{i,x}|\le\mu f_{i,z},|f_{i,y}|\le\mu f_{i,z},f_{i,z}\ge0$，共 $4\times5=20$ 个线性不等式，安全关键时换内逼近），ZMP 约束\emph{免费}满足。代价是变量从 $6$D 升到 $12$D；但\textbf{线性约束远比非线性约束快}，这个 trade-off 划算。

\begin{insight}[12 维力不是更精确，而是线性化参数化技巧]
\label{ins:fw-12d}
$12$ 维顶点力\emph{不是}更精确的力模型——它描述的物理现实与 $6$D wrench\emph{完全相同}（二者经 $\mathbf{w}=\mathbf{G}\mathbf{f}_{12}$ 互通）。它\emph{是}一种\textbf{把非线性约束（SOCP + ZMP 除法）变成线性约束（QP）的参数化技巧}。由于 $12-6=6$ 维力\emph{冗余}，须靠正则化在可行域内选偏好解：均匀正则 $\|\mathbf{f}_{12}\|^2$ 让法向力均布（ZMP 趋脚底中心、能效优先）；加权正则 $\sum_i w_i\|\mathbf{f}_i\|^2$ 引导 ZMP（如前掌权重高、利上坡防后倾，平衡优先）；最小法向力方差 $\sum_i(f_{i,z}-\bar f_z)^2$ 让压强最均（软地面保护优先）。无“最好”策略，取决于任务。
\end{insight}

% ════════════════════════════════════════════════════════════════
\section{力控任务的二端口统一视角}
\label{sec:fw-twoport}

把力控嵌入 WBC 有两种视角。\cref{sec:fw-contact6d} 是“力作为\emph{优化变量}”（TSID 把力放进 QP 代价/约束）；另一种是“力作为\emph{任务描述}”——声明“我想在此面上施 $20\,\mathrm{N}$ 法向力”，框架自动组装。后者以 mc\_rtc 的四种力控任务为代表，而它们的统一钥匙是\emph{二端口网络}（主体见 O 章 \cref{ch:operational-space}，本章只取归类骨架）。

\cref{tab:fw-fourtask} 把四种任务沿\emph{力控谱}（纯力跟踪 $\leftrightarrow$ 纯位置跟踪）排列，并点明它们都是\textbf{同一二端口网络在不同工作点的实例}：DampingTask 是阻抗的 $\mathbf{K}_d=0$ 特例（零刚度纯阻尼，碰撞后“刹车”）；AdmittanceTask 是导纳控制的 WBC 版（测力、出位置参考）；ImpedanceTask 是阻抗控制的 WBC 版（测运动、出力参考）；CoPTask 经 \cref{thm:fw-zmp} 的接触力分布约束实现 ZMP/CoP 跟踪。

\begin{table}[H]\centering\small
\caption{力控四任务的二端口归类（mc\_rtc，理论骨架）}
\label{tab:fw-fourtask}
\begin{tabular}{llll}
\toprule
任务 & 因果性 & 输入 $\to$ 输出 & 本质（二端口工作点） \\
\midrule
DampingTask & — & 速度 $\to$ 阻尼约束 & 阻抗 $\mathbf{K}_d=0$ 特例 \\
AdmittanceTask & 导纳 & 力参考 $\to$ 位置参考 & 导纳控制的 WBC 版 \\
ImpedanceTask & 阻抗 & 位姿+力参考 $\to$ 力参考 & 阻抗控制的 WBC 版 \\
CoPTask & — & CoP 目标 $\to$ 接触力分布 & ZMP/CoP 跟踪（\cref{thm:fw-zmp}） \\
\bottomrule
\end{tabular}
\end{table}

\begin{remark}[SE(3) 阻抗误差用本书右扰动对数映射]
\label{rmk:fw-se3-error}
ImpedanceTask 的位姿误差是 $\mathrm{SE}(3)$ 上的量，应用 \cref{ch:lie} 的对数映射而非逐元素相减。按 \cref{rmk:force-wbc-notation}，姿态误差取右扰动 $\mathbf{e}_R=\mathrm{Log}(\mathbf{R}^\top\mathbf{R}_d)$（机体系），恢复力矩 $+\mathbf{K}_R\mathbf{e}_R$；位置误差 $\mathbf{e}_x=\mathbf{x}_d-\mathbf{x}$。这样力控任务的几何与本书 $\SO(3)/\SE(3)$ 主线一致，避免源里 “current$-$desired vs desired$-$current $\pm$ 号”的约定混乱（源的多种约定都自洽，本书统一取一种）。
\end{remark}

% ════════════════════════════════════════════════════════════════
\section{阻抗控制与 WBC 的统一视角}
\label{sec:fw-unification}

本章的理论收口：证明\textbf{阻抗控制是 WBC 的单任务特例}，从而把上半章（机械臂阻抗）与下半章（浮基 WBC）、以及 \cref{ch:quad_wbc}（腿足 WBC）焊为一体。

\begin{theorem}[阻抗控制是 WBC-QP 的单任务退化]
\label{thm:fw-kkt}
考虑 WBC 加权 QP（决策变量 $\dot{\boldsymbol{\nu}},\boldsymbol{\tau},\boldsymbol{\lambda}$）
\begin{equation}
  \min_{\dot{\boldsymbol{\nu}},\boldsymbol{\tau},\boldsymbol{\lambda}}\ \|\mathbf{J}\dot{\boldsymbol{\nu}}+\dot{\mathbf{J}}\boldsymbol{\nu}-\ddot{\mathbf{x}}_{\rm ref}\|^2+w_\tau\|\boldsymbol{\tau}\|^2
  \quad \text{s.t.}\quad \mathbf{M}\dot{\boldsymbol{\nu}}+\mathbf{h}=\mathbf{S}^\top\boldsymbol{\tau}+\mathbf{J}_c^\top\boldsymbol{\lambda},
  \label{eq:fw-wbcqp}
\end{equation}
$\ddot{\mathbf{x}}_{\rm ref}=\ddot{\mathbf{x}}_d+\mathbf{K}_p\mathbf{e}_x+\mathbf{K}_d\dot{\mathbf{e}}_x$。在四条件下——\emph{(1)} 固定基座（$\mathbf{S}=\mathbf{I}$）、\emph{(2)} 无接触（$\boldsymbol{\lambda}=\mathbf{0}$）、\emph{(3)} 单任务、\emph{(4)} $w_\tau\to0$——其 KKT 解为
\begin{equation}
  \boldsymbol{\tau}=\mathbf{M}\mathbf{J}^+(\ddot{\mathbf{x}}_{\rm ref}-\dot{\mathbf{J}}\dot{\mathbf{q}})+\mathbf{h},
  \label{eq:fw-kkt-sol}
\end{equation}
与笛卡尔阻抗律具有相同的“PD 加速度参考 $+$ 逆动力学”结构。
\end{theorem}

\begin{proof}[KKT 推导]
四条件下问题退化为 $\min_{\ddot{\mathbf{q}},\boldsymbol{\tau}}\|\mathbf{J}\ddot{\mathbf{q}}+\dot{\mathbf{J}}\dot{\mathbf{q}}-\ddot{\mathbf{x}}_{\rm ref}\|^2$ s.t. $\mathbf{M}\ddot{\mathbf{q}}+\mathbf{h}=\boldsymbol{\tau}$。Lagrange 函数 $L=\|\mathbf{J}\ddot{\mathbf{q}}+\dot{\mathbf{J}}\dot{\mathbf{q}}-\ddot{\mathbf{x}}_{\rm ref}\|^2+\boldsymbol{\lambda}_L^\top(\mathbf{M}\ddot{\mathbf{q}}+\mathbf{h}-\boldsymbol{\tau})$。驻点条件：
\begin{align}
  \nabla_{\ddot{\mathbf{q}}}L&=2\mathbf{J}^\top(\mathbf{J}\ddot{\mathbf{q}}+\dot{\mathbf{J}}\dot{\mathbf{q}}-\ddot{\mathbf{x}}_{\rm ref})+\mathbf{M}^\top\boldsymbol{\lambda}_L=\mathbf{0},\\
  \nabla_{\boldsymbol{\tau}}L&=-\boldsymbol{\lambda}_L=\mathbf{0}\ \Rightarrow\ \boldsymbol{\lambda}_L=\mathbf{0}.
\end{align}
代回第一式得 $\mathbf{J}^\top(\mathbf{J}\ddot{\mathbf{q}}+\dot{\mathbf{J}}\dot{\mathbf{q}}-\ddot{\mathbf{x}}_{\rm ref})=\mathbf{0}$，即 $\ddot{\mathbf{q}}=\mathbf{J}^+(\ddot{\mathbf{x}}_{\rm ref}-\dot{\mathbf{J}}\dot{\mathbf{q}})$；代入动力学约束即 \cref{eq:fw-kkt-sol}。
\end{proof}

\begin{remark}[“同构”不等于“总等价”——伪逆度量的微妙]
\label{rmk:fw-not-equiv}
必须\emph{克制}地表述这一结论。\cref{eq:fw-kkt-sol} 用的是 Moore--Penrose 逆 $\mathbf{J}^+$（最小\emph{关节}范数解）；只有当 QP 代价采用\emph{动力学一致度量}时，才会出现 $\bar{\mathbf{J}}=\mathbf{M}^{-1}\mathbf{J}^\top\boldsymbol{\Lambda}$，进一步接近操作空间控制的 $\mathbf{J}^\top\boldsymbol{\Lambda}(\ddot{\mathbf{x}}_{\rm ref}-\dot{\mathbf{J}}\dot{\mathbf{q}})$（即真正的阻抗律）。此外真实 WBC 还受\emph{力矩正则化}（$w_\tau\neq0$）、\emph{约束活动集}、\emph{接触力变量}影响。故只能说二者\textbf{结构同构、特定度量下重合}，不能说“WBC $=$ 阻抗”。这一克制比绝对化说法更严谨。
\end{remark}

\begin{insight}[WBC 是阻抗控制在三个方向上的自然推广]
\label{ins:fw-generalization}
\cref{thm:fw-kkt} 揭示：WBC 不是“新发明”的力控方法，而是阻抗控制沿三个方向的自然推广——\emph{(1)} 固定基 $\to$ 浮基（引入 $\mathbf{S}^\top\boldsymbol{\tau}$ 与 $\mathbf{J}_c^\top\boldsymbol{\lambda}$）；\emph{(2)} 单任务 $\to$ 多任务（引入权重/优先级）；\emph{(3)} 无约束 $\to$ 有约束（引入 QP 与摩擦锥）。理解了这层推广关系，WBC 的每个组件“为什么在那里”就一目了然。反过来，二者是\emph{包含}而非\emph{替代}关系：固定基座单任务场景下，闭式阻抗比 WBC 简单、快、易调参（WBC 的每个任务本身可以是一个阻抗任务）。\cref{tab:fw-unify} 给出谱系两端的对照。
\end{insight}

\begin{table}[H]\centering\small
\caption{阻抗控制 $\to$ WBC 的统一谱系}
\label{tab:fw-unify}
\begin{tabular}{llll}
\toprule
特性 & 阻抗控制 & WBC 加权 QP & WBC HQP \\
\midrule
基座 & 固定 & 浮动 & 浮动 \\
任务数 & 1 & $N$（加权） & $N$（严格优先级） \\
接触力 & 忽略 & 优化变量 & 优化变量 \\
约束 & 无 & 摩擦锥 + 力矩限 & 摩擦锥 + 力矩限 \\
求解 & 解析公式 & $1$ 个 QP & $N$ 个 QP（字典序） \\
\bottomrule
\end{tabular}
\end{table}

\paragraph{建模即性能：消元如何决定能否实时。}
统一视角还藏着一个工程洞察。WBC-QP 的决策变量有两种取法：\emph{保留三元} $\mathbf{z}=(\dot{\boldsymbol{\nu}},\boldsymbol{\lambda},\boldsymbol{\tau})$（如 legged\_control，直观、可直接读出 $\boldsymbol{\tau}$），或\emph{消去} $\boldsymbol{\tau}$ 取 $\mathbf{z}=(\dot{\boldsymbol{\nu}},\boldsymbol{\lambda})$（如 TSID）。后者利用：一旦 $(\dot{\boldsymbol{\nu}},\boldsymbol{\lambda})$ 定，驱动力矩可由动力学\emph{唯一反推}
\begin{equation}
  \boldsymbol{\tau}=\mathbf{S}\big(\mathbf{M}\dot{\boldsymbol{\nu}}+\mathbf{h}-\mathbf{J}_c^\top\boldsymbol{\lambda}\big),
  \label{eq:fw-tau-elim}
\end{equation}
而基座的 $6$ 行（无驱动）转为对 $(\dot{\boldsymbol{\nu}},\boldsymbol{\lambda})$ 的等式约束。

\begin{insight}[同一物理问题，换数学表述影响工程生死]
\label{ins:fw-modeling}
以四足为例，保留三元约 $42$ 维、消元后约 $30$ 维。QP 求解时间大致随 $n^2\!\sim\!n^3$，$30$ vs $42$ 维可\textbf{快约 $30\%$--$50\%$}——在 $1\,\mathrm{kHz}$ 实时约束下，这往往就是“能不能跑起来”的差别。\textbf{好的数学建模直接影响工程性能}：同一物理问题换一种表述即降近 $30\%$ 维度。这不只是优雅，更是实时系统的生死线。HQP/NSP 字典序、级联零空间投影等机制本身见 \cref{ch:quad_wbc}，此处只取“阻抗是 WBC 单任务特例 + 消元洞察”这两点。
\end{insight}

% ════════════════════════════════════════════════════════════════
\section{本章脉络与展望}
\label{sec:fw-outlook}

\paragraph{一条主梁回望。}
本章把机器人力控织成一条连续的上升线（\cref{fig:fw-roadmap}）：纯位置控制在接触下“打架”$\to$ 力位混合\emph{分管}力与位置（需约束几何）$\to$ 主动刚度与并行控制免去硬切换却仍条件稳定 $\to$ 阻抗控制\emph{整形}力--运动关系并经无源性/LaSalle 拿到对全谱环境的渐近稳定 $\to$ 浮基升维引出 $12$ 维顶点力（ZMP 凸包隐式满足）与二端口四任务 $\to$ 最终经 KKT 证明\emph{阻抗是 WBC 的单任务特例}，与 \cref{ch:quad_wbc} 的腿足 WBC 在“多任务约束优化”处会合。

\begin{figure}[H]
\centering
\begin{tikzpicture}[font=\small, node distance=5mm and 9mm, >=Stealth,
  st/.style={draw, rounded corners, align=center, fill=main!7, draw=main!60!black, inner sep=4pt, font=\small},
  hi/.style={draw, rounded corners, align=center, fill=second!10, draw=second!60!black, inner sep=4pt, font=\small},
  ar/.style={->, thick, gray!60!black}]
  \node[st] (pos) {纯位置控制\\（接触“打架”）};
  \node[st, right=of pos] (hyb) {力位混合\\ 分管力/位};
  \node[st, right=of hyb] (imp) {阻抗控制\\ 整形力--运动};
  \node[hi, right=of imp] (wbc) {WBC\\ 多任务约束优化};
  \node[st, below=7mm of hyb] (stiff) {主动刚度/并行\\（条件稳定）};
  \node[st, below=7mm of imp] (pass) {无源性 + LaSalle\\（全谱稳定）};
  \node[hi, below=7mm of wbc] (c6d) {12 维顶点力\\ ZMP 凸包};
  \draw[ar] (pos)--(hyb); \draw[ar] (hyb)--(imp); \draw[ar] (imp)--(wbc);
  \draw[ar] (hyb.south)-- (stiff.north); \draw[ar] (imp.south)--(pass.north);
  \draw[ar] (wbc.south)--(c6d.north);
  \draw[ar] (pass.east)-- (c6d.west) node[midway, above, font=\footnotesize]{KKT 统一};
\end{tikzpicture}
\caption{本章主梁：从纯位置控制到 WBC 的力控复杂度谱。阻抗经无源性拿到全谱稳定，并经 KKT 与 WBC 在“多任务约束优化”处统一。}
\label{fig:fw-roadmap}
\end{figure}

\paragraph{前沿一瞥。}
WBC 与 MPC 的边界正在融合，几条代表线（仅点名，机制见各对应章）：\emph{VWBC}（value-function-based WBC）用 MPC 反向扫描得到的 $Q$ 函数（cost-to-go）替代手调权重、消除 WBC 调参，由 Cafe-MPC 的级联保真度框架提出（Li \& Wensing，arXiv 2024 / IEEE T-RO 2024\cite{li2024cafempc}）；\emph{全身 MPC} 直接用全身动力学做 MPC（借 MuJoCo 高效并行 + iLQR，2025）以求消除分层；\emph{Whole-Body MPPI}（Rapuano 2024）首次把基于采样的全身 MPC 部署到真实四足，天然适合 GPU 且对接触不连续鲁棒。这些指向控制部 MPC 深化 \cref{ch:mpc-advanced}、DDP/iLQR \cref{ch:ddp-ilqr} 与规划部路径积分 \cref{ch:mppi}。

\paragraph{承上启下。}
向上，本章用到的操作空间 $\boldsymbol{\Lambda}$/动力学一致逆/二端口/无源性二分法的\emph{通用理论}在 O 章 \cref{ch:operational-space}，摩擦锥/GIWC/CoP 完整理论在接触部 \cref{ch:contact-mechanics}，浮基动力学/LIPM/CMM 在 \cref{ch:reduced-models}，Lyapunov/LaSalle/Barbalat 工具在 \cref{ch:lyapunov-stability}，自适应/辨识在 \cref{ch:robust-hinf}，$\dot{\mathbf{M}}-2\mathbf{C}$ 反对称等动力学性质在 \cref{ch:geometric-mechanics}。向下，本章是机械臂力控的\emph{理论地基}：HQP/NSP/浮基 QP 的完整机制在 \cref{ch:quad_wbc}（腿足侧入口），复合系统（移动操作底盘 + 臂统一 WBC）的统一动力学与多模态 MPC 在 \cref{ch:unified-multimodal-mpc}。三章——本章（机械臂阻抗，最“软/简”端）、\cref{ch:quad_wbc}（腿足完整 WBC）、复合系统 \cref{ch:unified-multimodal-mpc}——共同构成一条贯通的\emph{力控复杂度谱}。

% ════════════════════════════════════════════════════════════════
\section*{本章小结}
\addcontentsline{toc}{section}{本章小结}

\begin{itemize}
  \item \textbf{力控的根本困难}：受约束方向上力与位置不能同时独立指定（\cref{sec:fw-history}）。五大经典算法是对此的五种回应。
  \item \textbf{力位混合}（Raibert--Craig，\cref{sec:fw-raibert}）：任务选择矩阵 $\mathbf{S}_\sigma$ 硬切换分管力/位；条件稳定、需约束几何；接口量纲二义性须先定后写（\cref{rmk:fw-interface}）。
  \item \textbf{主动刚度}（Salisbury，\cref{sec:fw-stiffness}）：纯虚拟弹簧 $\tilde{\mathbf{x}}=\mathbf{K}_x^{-1}\mathbf{f}_{\rm ext}$；无阻尼会弹跳，仅 Lyapunov 稳定——这是 Hogan 加阻尼的动机。\textbf{并行控制}（Chiaverini，\cref{sec:fw-parallel}）：全方向并行、靠力积分主导实现零稳态力误差、免约束几何。
  \item \textbf{阻抗控制}（Hogan，\cref{sec:fw-impedance,sec:fw-passivity}）：整形二阶力--运动关系（刚度是柔顺旋钮）；取 $\mathbf{M}_d=\boldsymbol{\Lambda}$ 免力传感；无源性 + LaSalle 给全局渐近稳定，Colgate--Hogan 给对全谱环境的耦合稳定。零空间须用动力学一致逆 $\bar{\mathbf{J}}$（\cref{rmk:fw-nullspace}）。
  \item \textbf{柔性关节与自适应}（\cref{sec:fw-flex-adaptive}）：Ott 的等效刚度硬上界 $K_{d,eq}\le K_s$；Slotine--Li 自适应靠线性参数化 + $\dot{\mathbf{M}}-2\mathbf{C}$ 反对称 + Barbalat，无 PE 仍跟踪收敛。
  \item \textbf{浮基升维}（\cref{sec:fw-why-wbc,sec:fw-contact6d,sec:fw-twoport,sec:fw-unification}）：欠驱动/多任务/接触力须优化三大困难 $\to$ WBC；$12$ 维顶点力把 ZMP 非线性约束化为 $f_{i,z}\ge0$（凸包隐式满足，\cref{thm:fw-zmp}）；四力控任务是二端口不同工作点；KKT 证明阻抗是 WBC 单任务特例（\cref{thm:fw-kkt}，但只“结构同构”非“总等价”）。
\end{itemize}

\section*{常见误解}
\addcontentsline{toc}{section}{常见误解}
\begin{enumerate}
  \item \textbf{“接触力太大就把位置环调软”}。仍用纯位置环只是降增益，受约束方向力与位置依旧耦合。应显式设阻抗（刚度 $\mathbf{K}_d^x$）或做混合/并行力位控。
  \item \textbf{“$\dot V\le0$ 就渐近稳定”}。只给 Lyapunov 稳定；须 LaSalle/Barbalat 升级（\cref{ins:fw-lasalle}）。忘给某旋转方向设阻尼 $\Rightarrow$ 该方向永振荡。
  \item \textbf{“零空间随便用 $\mathbf{J}^+$”}。Moore--Penrose 逆做零空间非动力学一致，会泄漏末端扰动；须用 $\bar{\mathbf{J}}=\mathbf{M}^{-1}\mathbf{J}^\top\boldsymbol{\Lambda}$（\cref{rmk:fw-nullspace}）。
  \item \textbf{“12 维力比 6D wrench 更精确”}。物理现实相同，$12$ 维只是把非线性约束线性化的参数化技巧（\cref{ins:fw-12d}）。
  \item \textbf{“WBC 比阻抗更好、可完全取代它”}。二者是\emph{包含}非\emph{替代}关系；固定基座单任务下闭式阻抗更简单、更快、更易调参，WBC 的每个任务本身可以是阻抗任务（\cref{ins:fw-generalization}）。
  \item \textbf{“可设的笛卡尔刚度无上限”}。柔性关节下 $K_{d,eq}\le K_s$ 是物理硬上界（\cref{thm:fw-ott}），强行超越致弹簧振荡与饱和。
  \item \textbf{“自适应控制一定收敛到真参数”}。无持续激励时参数未必收敛，但跟踪误差仍趋零（等效参数现象，\cref{ins:fw-pe}）。
\end{enumerate}

\section*{延伸阅读}
\addcontentsline{toc}{section}{延伸阅读}
\begin{itemize}
  \item \textbf{阻抗控制奠基}：Hogan《Impedance Control I/II/III》\cite{hogan1985impedance}——力--运动关系整形的原始论证；Colgate \& Hogan\cite{colgate1988robust} 的耦合稳定定理（无源 $\Rightarrow$ 稳定）。
  \item \textbf{力位混合与约束分析}：Raibert \& Craig\cite{raibert1981hybrid}（选择矩阵奠基）、Mason\cite{mason1981compliance}（自然/人工约束）、De Schutter \& Van Brussel\cite{deschutter1988compliant}（任务帧/柔顺帧）、Chiaverini \& Sciavicco\cite{chiaverini1993parallel}（并行力位）、Salisbury\cite{salisbury1980active}（主动刚度）。
  \item \textbf{操作空间与动力学一致逆}：Khatib\cite{khatib1987operational}——$\boldsymbol{\Lambda}=(\mathbf{J}\mathbf{M}^{-1}\mathbf{J}^\top)^{-1}$、唯一给末端零加速度的广义逆 $\bar{\mathbf{J}}$（本章笛卡尔阻抗与零空间的理论出处，详见 O 章 \cref{ch:operational-space}）。
  \item \textbf{柔性关节与自适应}：Ott《Cartesian Impedance Control of Redundant and Flexible-Joint Robots》（Springer STAR 49）\cite{ott2008cartesian} 与统一无源框架（T-RO 24(6)）\cite{ott2008unified}；Slotine \& Li\cite{slotine1987adaptive}（自适应控制 + $\dot{\mathbf{M}}-2\mathbf{C}$ 反对称）。
  \item \textbf{WBC/TSID/HQP}：TSID 优先级运动--力控（Del Prete 等，RAS 2015）\cite{delprete2015tsid}；多机器人/任务空间 QP 力控（对应 mc\_rtc 框架，Bouyarmane 等，IEEE T-RO 2019）\cite{bouyarmane2019qp}；多优先级零空间投影 WBC 概念化（Sentis \& Khatib，IJHR 2005）\cite{sentis2005synthesis}；HQP 快速在线生成（Escande 等，IJRR 2014）\cite{escande2014hierarchical}；腿足优化控制综述（Wensing 等，T-RO 2024）\cite{wensing2024legged}——HQP/NSP 完整机制见 \cref{ch:quad_wbc}。
  \item \textbf{碰撞检测与前沿}：De Luca 等残差/广义动量观测器碰撞检测（DLR-III，IROS 2006）\cite{deluca2006collision}（无传感碰撞检测，后扩展到浮动基座）；Mini Cheetah WBIC+MPC（Kim 等）\cite{kim2019wbic}、Momentum HQP（Herzog 等）\cite{herzog2016momentum}、Cafe-MPC/VWBC（Li \& Wensing）\cite{li2024cafempc}。
\end{itemize}

\section*{本章与后续章节的关系}
\addcontentsline{toc}{section}{本章与后续章节的关系}
\begin{table}[H]\centering\small
\begin{tabular}{lll}
\toprule
相关章节 & 关系 & 本章接口 \\
\midrule
\cref{ch:control} 控制导论 & 导论级阻抗/WBC 的完整深化 & \cref{eq:ctrl-impedance,eq:ctrl-osc,eq:ctrl-floatbase} \\
\cref{ch:quad_wbc} 腿足全身控制 & 腿足侧 WBC，HQP/NSP 机制复用 & \cref{sec:fw-why-wbc,sec:fw-unification} 互引 \\
\cref{ch:operational-space} 操作空间(O) & $\boldsymbol{\Lambda}$/$\bar{\mathbf{J}}$/二端口/无源性主体 & \cref{sec:fw-impedance-cart,sec:fw-twoport} \\
\cref{ch:contact-mechanics} 接触力学(N) & 摩擦锥/GIWC/CoP 完整理论 & \cref{sec:fw-why-wbc,sec:fw-contact6d} \\
\cref{ch:reduced-models} 浮基/简化模型(M) & 浮基动力学/LIPM/CMM 主体 & \cref{eq:fw-floatbase} \\
\cref{ch:lyapunov-stability} Lyapunov(D) & LaSalle/Barbalat 工具出处 & \cref{thm:fw-passivity,thm:fw-adaptive} \\
\cref{ch:unified-multimodal-mpc} 复合系统(Q) & 移动操作统一 WBC + 多模态 MPC & \cref{sec:fw-why-wbc} 三场景 \\
\bottomrule
\end{tabular}
\end{table}
